\documentclass{beamer}
\usepackage[utf8]{inputenc}
\usepackage[spanish]{babel}
\usepackage{listings}
\usepackage{xcolor}
\usepackage{tikz}
\usepackage{tcolorbox}
\usepackage{fontawesome5}

% Configuración de colores - Gruvbox Light con marrón principal
\definecolor{bg0}{HTML}{fbf1c7}
\definecolor{bg1}{HTML}{ebdbb2}
\definecolor{bg2}{HTML}{d5c4a1}
\definecolor{bg3}{HTML}{bdae93}
\definecolor{fg0}{HTML}{282828}
\definecolor{fg1}{HTML}{3c3836}
\definecolor{fg2}{HTML}{504945}
\definecolor{aqua}{HTML}{689d6a}
\definecolor{aquabright}{HTML}{427b58}
\definecolor{blue}{HTML}{458588}
\definecolor{red}{HTML}{cc241d}
\definecolor{orange}{HTML}{d65d0e}
\definecolor{yellow}{HTML}{d79921}
\definecolor{green}{HTML}{98971a}
\definecolor{purple}{HTML}{b16286}
\definecolor{brown}{HTML}{AB886D}
\definecolor{brownbright}{HTML}{8B6849}

% Colores para código
\definecolor{codegreen}{HTML}{98971a}
\definecolor{codegray}{HTML}{7c6f64}
\definecolor{codepurple}{HTML}{b16286}
\definecolor{backcolour}{HTML}{f9f5d7}

% Configuración de código
\lstdefinestyle{pythonstyle}{
    backgroundcolor=\color{backcolour},   
    commentstyle=\color{codegreen},
    keywordstyle=\color{brownbright}\bfseries,
    numberstyle=\tiny\color{codegray},
    stringstyle=\color{codepurple},
    basicstyle=\ttfamily\footnotesize\color{fg0},
    breakatwhitespace=false,         
    breaklines=true,                 
    captionpos=b,                    
    keepspaces=true,                 
    numbers=left,                    
    numbersep=3pt,                  
    showspaces=false,                
    showstringspaces=false,
    showtabs=false,                  
    tabsize=2,
    language=Python,
    frame=single,
    rulecolor=\color{brown},
    xleftmargin=8pt,
    xrightmargin=8pt,
    inputencoding=utf8,
    extendedchars=true,
    literate=
        {á}{{\'a}}1 {é}{{\'e}}1 {í}{{\'i}}1 {ó}{{\'o}}1 {ú}{{\'u}}1
        {Á}{{\'A}}1 {É}{{\'E}}1 {Í}{{\'I}}1 {Ó}{{\'O}}1 {Ú}{{\'U}}1
        {ñ}{{\~n}}1 {Ñ}{{\~N}}1
        {¿}{{?`}}1 {¡}{{!`}}1
}

\lstset{style=pythonstyle}

% Tema personalizado con marrón principal
\usetheme{Madrid}
\setbeamercolor{structure}{fg=brown}
\setbeamercolor{palette primary}{bg=brown,fg=bg0}
\setbeamercolor{palette secondary}{bg=brownbright,fg=bg0}
\setbeamercolor{palette tertiary}{bg=fg1,fg=bg0}
\setbeamercolor{palette quaternary}{bg=brown,fg=bg0}
\setbeamercolor{block title}{bg=brown,fg=bg0}
\setbeamercolor{block body}{bg=bg1,fg=fg0}
\setbeamercolor{frametitle}{bg=brown,fg=bg0}
\setbeamercolor{title}{fg=fg0}
\setbeamercolor{subtitle}{fg=fg1}
\setbeamercolor{normal text}{fg=fg0,bg=bg0}
\setbeamercolor{item}{fg=brown}
\setbeamercolor{section in toc}{fg=brown}

% Cajas personalizadas con marrón
\newtcolorbox{conceptbox}[1]{
    colback=bg1,
    colframe=brown,
    fonttitle=\bfseries\small,
    title=#1,
    coltitle=fg0,
    left=3pt,
    right=3pt,
    top=3pt,
    bottom=3pt,
    boxsep=2pt
}

\newtcolorbox{notebox}{
    colback=yellow!15,
    colframe=orange,
    leftrule=2mm,
    rightrule=0mm,
    toprule=0mm,
    bottomrule=0mm,
    arc=0mm,
    left=3pt,
    right=3pt,
    top=2pt,
    bottom=2pt
}

\title{\textbf{Introducción a la Programación con Python}}
\subtitle{Semana 0: Conceptos Fundamentales}
\author{$\nabla$Nablabs}
\date{}
\begin{document}

\begin{frame}
\titlepage
\end{frame}

\begin{frame}
\frametitle{Contenido}
\tableofcontents
\end{frame}

\section{¿Qué es Python?}

\begin{frame}{Python: Un lenguaje amigable}
\begin{conceptbox}{¿Por qué Python?}
\small
\begin{itemize}
    \item[\faCheckCircle] Lenguaje de programación de alto nivel
    \item[\faCheckCircle] Sintaxis clara y legible
    \item[\faCheckCircle] Muy versátil: web, ciencia de datos, IA
    \item[\faCheckCircle] Gran comunidad y documentación
\end{itemize}
\end{conceptbox}

\vspace{0.2cm}
\begin{center}
\large\textcolor{brown}{\faCode \ Python es ideal para principiantes}
\end{center}
\end{frame}

\section{Funciones Básicas}

\begin{frame}[fragile]{La función print() \ \faTerminal}

\begin{conceptbox}{¿Qué hace?}
Muestra información en pantalla
\end{conceptbox}

\vspace{-0.2cm}
\begin{lstlisting}
print("Hola mundo")
print(42)
print("Python", "es", "genial")
\end{lstlisting}

\vspace{-0.2cm}
\begin{tcolorbox}[colback=green!10,colframe=green,title={\small\faArrowCircleRight \ Salida},coltitle=fg0,left=3pt,right=3pt,top=2pt,bottom=2pt,boxsep=1pt]
\small\texttt{Hola mundo\\
42\\
Python es genial}
\end{tcolorbox}
\end{frame}

\begin{frame}[fragile]{La función input() \ \faKeyboard}

\begin{conceptbox}{¿Qué hace?}
Lee datos ingresados por el usuario
\end{conceptbox}

\vspace{-0.2cm}
\begin{lstlisting}
nombre = input("¿Cómo te llamas? ")
print("Hola", nombre)
\end{lstlisting}

\vspace{0.2cm}
\begin{notebox}
\small\faExclamationTriangle \ \textbf{Nota:} \texttt{input()} siempre retorna un string
\end{notebox}
\end{frame}

\section{Variables}

\begin{frame}[fragile]{Variables: Guardando información \ \faDatabase}

\begin{conceptbox}{¿Qué son las variables?}
Contenedores que almacenan datos en la memoria
\end{conceptbox}

\vspace{-0.2cm}
\begin{lstlisting}
nombre = "Ana"
edad = 20
altura = 1.65
estudiante = True
\end{lstlisting}

\vspace{-0.1cm}
\begin{block}{Reglas para nombres de variables}
\small
\begin{itemize}
    \item[\faTimes] No pueden empezar con números
    \item[\faCheck] Usan letras, números y guiones bajos
    \item[\faExclamation] Sensibles a mayúsculas: \texttt{nombre} $\neq$ \texttt{Nombre}
\end{itemize}
\end{block}
\end{frame}

\section{Tipos de Datos}

\begin{frame}[fragile]{Tipos de datos básicos \ \faCubes}

\begin{columns}[T]
\column{0.48\textwidth}
\begin{conceptbox}{\faFont \ String (str)}
\vspace{-0.15cm}
\begin{lstlisting}
mensaje = "Hola"
\end{lstlisting}
\end{conceptbox}

\vspace{0.1cm}
\begin{conceptbox}{\faHashtag \ Integer (int)}
\vspace{-0.15cm}
\begin{lstlisting}
edad = 25
\end{lstlisting}
\end{conceptbox}

\column{0.48\textwidth}
\begin{conceptbox}{\faCalculator \ Float}
\vspace{-0.15cm}
\begin{lstlisting}
precio = 19.99
\end{lstlisting}
\end{conceptbox}

\vspace{0.1cm}
\begin{conceptbox}{\faToggleOn \ Boolean}
\vspace{-0.15cm}
\begin{lstlisting}
activo = True
\end{lstlisting}
\end{conceptbox}
\end{columns}
\end{frame}

\begin{frame}[fragile]{Conversión de tipos \ \faExchangeAlt}

\begin{conceptbox}{¿Para qué convertir tipos?}
Para realizar operaciones entre tipos compatibles
\end{conceptbox}

\vspace{-0.2cm}
\begin{lstlisting}
# String a entero
edad = int("25")  # edad = 25

# String a float
precio = float("19.99")  # 19.99

# Número a string
numero = str(100)  # "100"
\end{lstlisting}

\vspace{-0.1cm}
\begin{notebox}
\small\faLightbulb \ \textbf{Útil con input():} \texttt{edad = int(input("Edad: "))}
\end{notebox}
\end{frame}

\section{Comentarios}

\begin{frame}[fragile]{Comentarios en el código \ \faComment}

\begin{conceptbox}{¿Para qué sirven?}
Explican qué hace el código (Python los ignora)
\end{conceptbox}

\vspace{-0.2cm}
\begin{lstlisting}
# Comentario de una línea
nombre = "Carlos"  # Al final de línea

"""
Comentario
de múltiples líneas
"""
\end{lstlisting}

\vspace{-0.1cm}
\begin{block}{\faQuestionCircle \ ¿Por qué comentar?}
\small
\begin{itemize}
    \item Documenta tu código
    \item Ayuda a otros programadores
    \item Te ayuda a recordar tu lógica
\end{itemize}
\end{block}
\end{frame}

\section{Formateo de Strings}

\begin{frame}[fragile]{F-strings: Formato moderno \ \faMagic}

\begin{conceptbox}{Ventaja}
Insertar variables en strings de forma elegante
\end{conceptbox}

\vspace{-0.2cm}
\begin{lstlisting}
nombre = "María"
edad = 22

# Método antiguo (NO recomendado)
mensaje = "Hola " + nombre

# F-string (RECOMENDADO)
mensaje = f"Hola {nombre}, {edad} años"
print(mensaje)
\end{lstlisting}

\vspace{-0.2cm}
\begin{tcolorbox}[colback=green!10,colframe=green,title={\small\faArrowCircleRight \ Salida},coltitle=fg0,left=3pt,right=3pt,top=2pt,bottom=2pt]
\small\texttt{Hola María, 22 años}
\end{tcolorbox}
\end{frame}

\section{Métodos de Strings}

\begin{frame}[fragile]{Métodos útiles de strings \ \faTools}

\begin{lstlisting}
texto = "  hola mundo  "

# Eliminar espacios
limpio = texto.strip()  # "hola mundo"

# Primera mayúscula de cada palabra
titulo = texto.title()  # "Hola Mundo"

# Solo primera mayúscula
capital = texto.capitalize()  # "Hola mundo"

# Dividir en lista
palabras = texto.split()  # ["hola", "mundo"]
\end{lstlisting}
\end{frame}

\begin{frame}[fragile]{Ejemplo práctico \ \faUserEdit}

\begin{lstlisting}
nombre = input("Tu nombre: ")

# Limpiar y formatear
nombre_limpio = nombre.strip().title()

print(f"Bienvenido/a, {nombre_limpio}!")
\end{lstlisting}

\vspace{0.2cm}
\begin{columns}[T]
\column{0.48\textwidth}
\begin{tcolorbox}[colback=blue!10,colframe=blue,title={\small\faArrowDown \ Entrada},coltitle=fg0,left=2pt,right=2pt,top=2pt,bottom=2pt]
\small\texttt{\ \ juan pérez\ \ }
\end{tcolorbox}

\column{0.48\textwidth}
\begin{tcolorbox}[colback=green!10,colframe=green,title={\small\faArrowUp \ Salida},coltitle=fg0,left=2pt,right=2pt,top=2pt,bottom=2pt]
\small\texttt{Bienvenido/a, Juan Pérez!}
\end{tcolorbox}
\end{columns}
\end{frame}

\section{Operadores Aritméticos}

\begin{frame}[fragile]{Operaciones matemáticas \ \faCalculator}

\begin{lstlisting}
a = 10
b = 3

suma = a + b        # 13
resta = a - b       # 7
multiplicacion = a * b  # 30
division = a / b    # 3.333...
division_entera = a // b  # 3
modulo = a % b      # 1 (residuo)
potencia = a ** b   # 1000
\end{lstlisting}

\vspace{-0.1cm}
\begin{notebox}
\small\faInfoCircle \ El operador \texttt{\%} devuelve el residuo de la división
\end{notebox}
\end{frame}

\begin{frame}[fragile]{Ejemplo: Calculadora simple \ \faDesktop}

\begin{lstlisting}
# Solicitar números
num1 = float(input("Primer número: "))
num2 = float(input("Segundo número: "))

# Calcular operaciones
suma = num1 + num2
producto = num1 * num2

# Mostrar resultados
print(f"Suma: {suma}")
print(f"Producto: {producto}")
\end{lstlisting}

\vspace{-0.1cm}
\begin{center}
\textcolor{brown}{\faRocket \ ¡Inténtalo tú mismo!}
\end{center}
\end{frame}

\section{Ejemplo Integrador}

\begin{frame}[fragile]{Programa completo \ \faIdCard}

\begin{lstlisting}
# Solicitar información
nombre = input("Nombre: ").strip().title()
edad = int(input("Edad: "))
ciudad = input("Ciudad: ").strip().title()

# Calcular año de nacimiento
anio_actual = 2025
anio_nacimiento = anio_actual - edad

# Mostrar tarjeta
print("\n--- TARJETA ---")
print(f"Nombre: {nombre}")
print(f"Edad: {edad} años")
print(f"Ciudad: {ciudad}")
print(f"Nacimiento: {anio_nacimiento}")
\end{lstlisting}
\end{frame}

\begin{frame}{Resumen \ \faListUl}

\begin{conceptbox}{¿Qué aprendimos hoy?}
\small
\begin{itemize}
    \item[\faCheck] Funciones: \texttt{print()}, \texttt{input()}
    \item[\faCheck] Variables y tipos de datos
    \item[\faCheck] Conversión de tipos (casting)
    \item[\faCheck] Comentarios
    \item[\faCheck] F-strings para formateo
    \item[\faCheck] Métodos de strings
    \item[\faCheck] Operadores aritméticos
\end{itemize}
\end{conceptbox}

\vspace{0.3cm}
\begin{center}
\large\textcolor{brown}{\faCode \ ¡Siguiente paso: Practicar!}
\end{center}
\end{frame}

\begin{frame}
\begin{center}
\Huge\textcolor{brown}{\faQuestionCircle}\\
\vspace{0.5cm}
\Huge ¡Preguntas!\\
\vspace{1cm}
\Large\textcolor{brownbright}{Gracias por su atención}
\end{center}
\end{frame}

\end{document}